\section{Methodology}
\label{sec:method}

In this section, we present the formal mathematical framework of \textbf{Quantum Wavefunction Superposition Merging (QWS-Merge)}. We transition from static weight compromises to dynamic parameter superposition.

\subsection{Problem Formulation}
Let $W_{base}^{(l)}$ denote the pre-trained base network weights at layer $l \in \{1, \dots, L\}$. Suppose we have fine-tuned $K$ independent task-specific expert models starting from the same base model. For each task $k \in \{1, \dots, K\}$, the task-specific weights at layer $l$ are denoted as $W_k^{(l)}$. Following Task Arithmetic \cite{ilharco2022editing}, we define the task vector $V_k^{(l)}$ as the parameter difference between the $k$-th expert and the base network:
\begin{equation}
    V_k^{(l)} = W_k^{(l)} - W_{base}^{(l)}.
\end{equation}
Rather than combining these task vectors using static coefficients, QWS-Merge defines dynamic, input-dependent merging coefficients determined via wave-like phase interference.

\subsection{Input Phase State Extraction}
Let $x \in \mathbb{R}^{B \times C \times H \times W}$ be a batch of input images. The input is first processed by the backbone's frozen patch embedding layer to extract patch tokens:
\begin{equation}
    H_0 = \text{PatchEmbed}(x) \in \mathbb{R}^{B \times N \times D},
\end{equation}
where $N$ is the number of patch tokens and $D$ is the embedding dimension. We extract a global representation $z(x) \in \mathbb{R}^{B \times D}$ by computing the spatial average over the patch tokens:
\begin{equation}
    z(x)_b = \frac{1}{N} \sum_{n=1}^N H_{0, b, n, :} \quad \forall b \in \{1, \dots, B\}.
\end{equation}
We then project this global representation into a low-dimensional phase-state space ($d$-dimensional, with $d = K = 4$) via a frozen random projection matrix $P \in \mathbb{R}^{D \times d}$ and normalize it to the unit sphere:
\begin{equation}
    \tilde{\psi}(x) = z(x) P \in \mathbb{R}^{B \times d},
\end{equation}
\begin{equation}
    \psi(x)_b = \frac{\tilde{\psi}(x)_b}{\|\tilde{\psi}(x)_b\|_2 + \epsilon} \quad \forall b \in \{1, \dots, B\},
\end{equation}
where $\epsilon = 10^{-8}$ is a small numerical stabilizer. The vector $\psi(x)_b$ represents the input phase state of the $b$-th batch sample.

\subsection{Quantum Phase-Coherent Overlap}
For each layer group $l \in \{1, \dots, L\}$ and task expert $k \in \{1, \dots, K\}$, we define a trainable phase-basis vector $\Phi_k^{(l)} \in \mathbb{R}^d$ and normalize it to the unit sphere:
\begin{equation}
    \hat{\Phi}_k^{(l)} = \frac{\Phi_k^{(l)}}{\|\Phi_k^{(l)}\|_2 + \epsilon} \in \mathbb{R}^d.
\end{equation}
The dynamic probability amplitude (sample-level coefficient) $\alpha_{k, b}(l)$ representing the contribution of task $k$'s expertise at layer $l$ for sample $b$ is computed as:
\begin{equation}
    \alpha_{k, b}(l) = R_k^{(l)} \cos\left( \omega \langle \psi(x)_b, \hat{\Phi}_k^{(l)} \rangle + \phi_k^{(l)} \right),
\end{equation}
where:
\begin{itemize}
    \item $R_k^{(l)} \in \mathbb{R}$ is a learned scaling amplitude for task $k$ at layer $l$, initialized to $0.3$.
    \item $\omega = \pi$ is a fixed wave frequency scaling factor.
    \item $\phi_k^{(l)} \in \mathbb{R}$ is a learned task-specific phase offset bias for task $k$ at layer $l$, initialized to $0.0$.
    \item $\langle \cdot, \cdot \rangle$ denotes the standard inner product in Euclidean space.
\end{itemize}
This cosine formulation simulates wave-like interference. When the input phase state $\psi(x)_b$ is highly aligned with the layer's task phase basis $\hat{\Phi}_k^{(l)}$, constructive interference occurs, yielding a high coefficient amplitude. Conversely, destructive interference reduces the coefficient, preventing conflicting expert weights from corrupting the activation pathway.

\subsection{Wavefunction Collapse \& Weight Measurement}
To resolve the computational challenge of sample-wise batch processing on standard deep learning accelerators, we perform a mean-measurement across the batch. This collapse projects the superposed wave state into a single batch-level classical coefficient $\bar{\alpha}_k(l)$:
\begin{equation}
    \bar{\alpha}_k(l) = \frac{1}{B} \sum_{b=1}^B \alpha_{k, b}(l).
\end{equation}
The dynamically assembled classical weight matrix for layer $l$ used to process the batch $x$ is given by:
\begin{equation}
    W_{merged}^{(l)}(x) = W_{base}^{(l)} + \sum_{k=1}^K \bar{\alpha}_k(l) V_k^{(l)}.
\end{equation}
We then proceed with the standard forward pass of the $l$-th layer block using $W_{merged}^{(l)}(x)$. Since the parameter assembly requires only scalar-tensor additions and multiplications, the computational overhead is negligible compared to the forward pass itself.

\subsection{Classical Routing Baseline: Linear Router}
To benchmark the necessity of the quantum-inspired wave phase formulation, we introduce a classical soft-routing baseline, the Linear Router. The Linear Router maps the global pooled representation $z(x) \in \mathbb{R}^{B \times D}$ directly to the task coefficient space via a standard linear layer:
\begin{equation}
    o(x) = z(x) W_{route} + b_{route} \in \mathbb{R}^{B \times K},
\end{equation}
where $W_{route} \in \mathbb{R}^{D \times K}$ and $b_{route} \in \mathbb{R}^K$ are trainable parameters. The sample-level soft routing coefficients are obtained via Softmax:
\begin{equation}
    \alpha_{k, b} = \text{Softmax}(o(x)_b)_k.
\end{equation}
Like QWS-Merge, these sample-level coefficients are averaged across the batch to collapse to a batch-level weight representation:
\begin{equation}
    \bar{\alpha}_k = \frac{1}{B} \sum_{b=1}^B \alpha_{k, b}.
\end{equation}
While simple and intuitive, the classical Linear Router suffers from two major design limitations:
\begin{enumerate}
    \item \textbf{High Overfitting Risk in Data-Scarce Regimes:} The projection matrix $W_{route}$ is completely unconstrained, with $D \times K = 192 \times 4 = 768$ parameters. Under low-data calibration, this unconstrained parameter space is highly susceptible to overfitting, leading to representation collapse under out-of-distribution or high-conflict tasks.
    \item \textbf{No Layer-Wise Specialization:} The collapsed routing coefficients $\bar{\alpha}_k$ are global and applied identically across all $L$ layers of the backbone. This denies the network the ability to form layer-wise, localized functional specialization.
\end{enumerate}
By contrast, QWS-Merge leverages layer-wise trainable phase-basis vectors to support fine-grained, localized routing while bounding the routing mechanism inside a highly regularized, spherical wave phase-interference subspace.

\subsection{Parameter and Optimization Efficiency}
Across $L=14$ layer groups and $K=4$ task experts, the trainable parameter set of QWS-Merge consists of:
\begin{itemize}
    \item Amplitudes $R_k^{(l)}$: $14 \times 4 = 56$ parameters.
    \item Phases $\Phi_k^{(l)}$: $14 \times 4 \times 4 = 224$ parameters.
    \item Biases $\phi_k^{(l)}$: $14 \times 4 = 56$ parameters.
\end{itemize}
This results in exactly $336$ total trainable parameters, which is less than half of the $772$ parameters of the Linear Router. This ultra-compact parameter footprint restricts the optimization to a highly specialized manifold of parameter states. Consequently, QWS-Merge can be successfully optimized using standard Adam on a tiny validation set (16 samples per task, 64 total) within 100 steps, running in under 30 seconds on a standard processor. This avoids the Overfitting-Optimizer Paradox that affects high-dimensional parameter search spaces under data-scarce regimes.
